\section{Summary}\label{sec:tsps-summary}

This chapter addresses the gap between structure preserving signatures and their
threshold counterparts. Structure preserving signatures integrate naturally with
non-interactive zero knowledge proofs, which makes them a natural building block
for privacy preserving applications such as group signatures, anonymous and
delegatable credentials, and privacy preserving payment systems. Many of these
applications place a certifier, endorser or log provider at the point of
signing. That role is not one the system can do without, but the capability
attached to it can be narrowed. Distributing the signing key across several
parties leaves the role intact, while ensuring that none of them can sign alone
or ever holds the key. Efficient threshold structure preserving signatures had,
until now, remained an elusive goal.

We gave two threshold structure preserving signature schemes, built from the
structure preserving signatures of Abe, Groth, Haralambiev, and Ohkubo and of
Kiltz, Pan, and Wee, together with generic MPC interfaces over both field and
group elements~\cite{C:AGHO11,C:KilPanWee15}. Both constructions push the bulk
of their computation into an offline pre-processing phase, ahead of the message
to be signed becoming known, so that their online phase remains non-interactive
and comparable in cost to ordinary, non-threshold signing. We proved that both
constructions are correct, that no set of at most $t$ signers learns anything
about the signing key, that thresholdisation leaves the public key distribution
and the verification algorithm unchanged, and that both achieve adaptive TS-UF-1
security, inheriting the unforgeability of their underlying structure preserving
signatures whenever the MPC interfaces they build on are securely realised.

We evaluated both schemes using two realisations of these interfaces, ATLAS in
the honest majority setting and MASCOT in the dishonest majority setting, and
compared them against existing threshold structure preserving
signatures~\cite{C:GLOPS21,CCS:KelOrsSch16}. Our scheme built from the signature
of Abe, Groth, Haralambiev, and Ohkubo achieves optimally small threshold
signatures of three group elements, matching the size of its non-threshold
counterpart. Our scheme built from that of Kiltz, Pan, and Wee instead relies
only on the standard SXDH assumption~\cite{C:AGHO11,C:KilPanWee15}. In both
cases, the offline phase scales with the number of parties, while the online
phase remains constant-time and non-interactive, making the schemes practical
for high-throughput applications such as payments and transparency logs that
experience bursty traffic.

Beyond the constructions themselves, we set out the multiparty computation our
schemes consume, over field elements and over group elements alike, and extended
MP-SPDZ so that the group element interfaces are actually available to run them.
Computing on shared group elements is well studied and we claim no novelty for
it, as \Cref{sec:tsps-mpc} explains. What the chapter contributes is the
observation that structure preserving signing decomposes so that its
message-dependent part is entirely local, which is what makes the online phase
free of interaction, together with the two schemes that exploit
it~\cite{FOCS:Feldman87,IMA:SmaAla19,LATIN:ADEO21,NDSS:GPVRNC23}. We also
implemented previously proposed threshold structure preserving signature schemes
to allow for a concrete efficiency comparison.

The chapter also illustrates the wider argument of this thesis. The certifiers
these applications rely on, log providers, endorsers, ledger administrators,
were not removed. Their authority was divided, so that no one of them can
certify alone, and the structure preserving property kept the values they
certify out of their reach, so that they hold authority without knowledge.  What
made the arrangement deployable was the second constraint we imposed on
ourselves, that the online phase require no interaction at all: an institution
can only be given a role it can afford to perform.

Extending our approach to other structure preserving signatures, and to the
applications discussed in \Cref{sec:tsps-introduction}, would be valuable
directions for future work. Two gaps also remain between the analysis and the
implementation. The first is a realisation of the MPC interfaces secure against
adaptive corruption, which would let the concrete schemes inherit the adaptive
notion rather than its static analogue. The second is a reduction avoiding the
combinatorial loss of \Cref{thm:tsps-construction-security}, whether through an
equivocable sharing or an assumption of secure erasures.
